% Chapter 10: Temporal Transport and Persistence
% Part III: Multi-Formation and Temporal Theory

\chapter{시간적 전달과 지속성}
\label{ch:temporal-transport}

\begin{quote}
\textit{``동일한 것이 다시 나타난 것인가, 아니면 비슷한 것이 새로 생겨난 것인가?''}
\end{quote}

\vspace{0.5em}
제9장에서 우리는 복수의 형성이 공간적으로 공존하는 구조를 다루었다. 이 장에서는 더 근본적인 질문으로 나아간다: 형성은 시간을 가로질러 어떻게 \emph{지속}하는가?

전통적 추적(tracking) 접근은 각 개체에 고유 식별자를 부여하고 시간 프레임 간에 대응시킨다. 그러나 SCC는 이산적 객체를 전제하지 않으므로, 추적 ID를 사용할 수 없다. 대신, 형성의 \textbf{응집적 구조 자체}가 시간적 정체성의 근거가 된다. 이것이 \textbf{시간적 전달 핵(temporal transport kernel)} $\mathbf{M}_{t \to s}$의 역할이다.


\section{시간적 정체성의 문제}
\label{sec:temporal-identity-problem}

\subsection{추적 ID 없는 정체성}

전통적 컴퓨터 비전에서 ``같은 물건''의 시간적 정체성은 라벨(label)에 의존한다: $\mathrm{ID}_t = \mathrm{ID}_{t+1}$이면 같은 물건이다. 그러나 이 접근에는 세 가지 근본적 문제가 있다:

\begin{enumerate}
    \item \textbf{순환성}: 추적은 ``같은 물건''이라는 개념을 전제하지만, 이 개념 자체가 추적이 답해야 할 질문이다.
    \item \textbf{분열과 합병}: 물건이 분열하면 ID는 어느 쪽에 귀속되는가? 합병하면 어떤 ID가 살아남는가?
    \item \textbf{점진적 변형}: 테세우스의 배---구성 요소가 하나씩 교체되면 언제 ``다른 물건''이 되는가?
\end{enumerate}

SCC는 이러한 문제를 회피한다. 정체성은 라벨이 아니라 \textbf{구조적 계승(structural inheritance)}에 의해 정의된다.

\subsection{발생적 정체성 (Genidentity)}

SCC의 시간적 정체성 개념은 철학적 \emph{발생적 정체성(genidentity)} 개념과 깊이 연결된다. 라이헨바흐(Reichenbach)가 제안한 이 개념에 따르면, 동일성은 ``같은 것이 다시 나타남''이 아니라 ``관계적 구조가 계승됨''에 의해 성립한다.

형식적으로, 시각 $t$의 형성 $u_t$와 시각 $s > t$의 형성 $u_s$ 사이의 정체성은 다음에 의해 성립한다:
\begin{center}
\emph{$u_t$의 응집적 핵심(cohesive core)이 전달 핵 $\mathbf{M}_{t \to s}$를 통해\\
$u_s$의 핵심으로 ``계승(inherited)''될 때, 형성은 시간적으로 \textbf{지속(persist)}한다.}
\end{center}

\subsection{왜 공간적 근접만으로는 부족한가}

직관적으로, ``가까이 있으면 같은 형성''이라고 할 수 있지 않을까? 두 가지 반례가 이를 부정한다:

\begin{enumerate}
    \item \textbf{교차}: 두 형성이 서로 지나치는 경우, 공간적 근접은 잘못된 대응을 만든다.
    \item \textbf{구조 변화}: 형성이 제자리에서 구조적으로 변형되는 경우(예: 핵심이 경계로 전환), 공간적 위치는 동일하나 구조적 계승은 단절된다.
\end{enumerate}

따라서 시간적 정체성은 공간적 근접이 아닌 \textbf{응집적 구조의 유사성}에 기반해야 한다. 이것이 응집 지문(cohesion fingerprint)의 동기이다.

\begin{example}[교차 문제의 구체적 예시]
\label{ex:crossing-problem}
$20 \times 20$ 격자에서 두 형성 $u^1, u^2$가 각각 좌측과 우측에서 출발하여 서로를 향해 이동한다고 하자.

\begin{itemize}
    \item 시각 $t_0$: 형성 1의 핵심 = $(5, 10)$, 형성 2의 핵심 = $(15, 10)$
    \item 시각 $t_1$: 형성 1 → $(10, 8)$, 형성 2 → $(10, 12)$ (교차 중)
    \item 시각 $t_2$: 형성 1 → $(15, 10)$, 형성 2 → $(5, 10)$ (교차 완료)
\end{itemize}

공간 근접 기반 대응은 $t_1$에서 두 형성이 가까우므로 혼동을 일으킨다. 최악의 경우, 형성 1의 $t_0$ 상태를 형성 2의 $t_2$ 상태에 대응시키는 ``정체성 교환''이 발생한다.

응집 지문 기반 전달은 이 문제를 해결한다: 각 형성의 닫힘 패턴($\mathrm{Cl}(u^k)$)과 구별 패턴($D_k$)이 형성의 ``구조적 서명''으로 작용하여, 공간적으로 가깝더라도 구조적으로 다른 형성을 올바르게 대응시킨다.
\end{example}


\section{응집 지문과 자기-참조적 비용}
\label{sec:cohesion-fingerprint}

\subsection{3성분 응집 지문}

각 지점 $x \in X_t$에 대해, 그 응집적 상태를 3차원 벡터로 요약한다:
\begin{equation}
    \varphi_t(x) = \Big(\, u_t(x),\; \mathrm{Cl}_t(u_t)(x),\; D_t(x;\, 1 - u_t) \,\Big) \in [0,1]^3
    \label{eq:fingerprint}
\end{equation}

세 성분의 의미는 다음과 같다:
\begin{itemize}
    \item $u_t(x)$: 해당 지점의 원시적(raw) 응집 강도
    \item $\mathrm{Cl}_t(u_t)(x)$: 이웃으로부터의 관계적 지지를 반영한 완성된 응집
    \item $D_t(x;\, 1 - u_t)$: 외부와의 구조적 비대칭---내부와 외부의 구별도
\end{itemize}

세 성분은 각각 독립적인 구조적 차원을 포착한다. 단순히 $u_t(x)$만 사용하면 원시적 강도만 비교하게 되나, 닫힘과 구별을 포함함으로써 형성의 \emph{관계적 구조}를 비교할 수 있다.

\subsection{왜 공동소속 $C(x,x)$는 제외되었나}

초기 설계에서는 4성분 지문이 고려되었으나, 실험적 분석 결과 공동소속 대각선 $C_t(x,x)$는 다음과 같은 이유로 제외되었다:
\begin{itemize}
    \item \textbf{기여도 $< 0.4\%$}: 전체 전달 비용에서 $C(x,x)$의 기여가 미미
    \item \textbf{수치 불안정}: 야코비안 노름 $\approx 9300$으로, 미소 변동에 극도로 민감
    \item \textbf{이론적 격하}: v2.1에서 $C_t$가 형식 우주에서 파생적 진단으로 격하
\end{itemize}

3성분 지문은 이론적 간결성과 수치적 안정성을 모두 만족한다.

\begin{example}[응집 지문의 수치적 계산]
\label{ex:fingerprint-numerical}
$10 \times 10$ 격자에서 최적화된 단일 형성의 여러 지점에서 응집 지문을 계산한다:

\begin{center}
\begin{tabular}{lcccl}
\hline
\textbf{지점 유형} & $u(x)$ & $\mathrm{Cl}(u)(x)$ & $D(x; 1\!-\!u)$ & $\varphi(x)$ \\
\hline
깊은 핵심 & 0.97 & 0.96 & 0.91 & $(0.97, 0.96, 0.91)$ \\
얕은 핵심 & 0.85 & 0.88 & 0.72 & $(0.85, 0.88, 0.72)$ \\
경계 & 0.45 & 0.51 & 0.31 & $(0.45, 0.51, 0.31)$ \\
외부 인접 & 0.08 & 0.12 & 0.05 & $(0.08, 0.12, 0.05)$ \\
먼 외부 & 0.01 & 0.02 & 0.01 & $(0.01, 0.02, 0.01)$ \\
\hline
\end{tabular}
\end{center}

주목할 점: 깊은 핵심의 지문 $(0.97, 0.96, 0.91)$과 경계의 지문 $(0.45, 0.51, 0.31)$ 사이의 제곱 거리는:
\begin{equation}
    \|\varphi_{\mathrm{core}} - \varphi_{\mathrm{bd}}\|^2 = (0.52)^2 + (0.45)^2 + (0.60)^2 = 0.833 \nonumber
\end{equation}
반면, 두 시각에서의 깊은 핵심 사이의 거리(형성이 약간 이동한 경우)는:
\begin{equation}
    \|\varphi^t_{\mathrm{core}} - \varphi^s_{\mathrm{core}}\|^2 \approx (0.02)^2 + (0.01)^2 + (0.03)^2 = 0.0014 \nonumber
\end{equation}
이 600배의 비용 차이가 전달을 핵심-핵심으로 집중시키는 근거이다.
\end{example}

\subsection{자기-참조적 비용}

지문을 이용한 전달 비용은 다음과 같이 정의된다:
\begin{equation}
    c(x, y) = \|\varphi_t(x) - \varphi_s(y)\|^2 = \sum_{i=1}^{3} \big(\varphi_t^{(i)}(x) - \varphi_s^{(i)}(y)\big)^2
    \label{eq:transport-cost}
\end{equation}

\textbf{핵심적 자기-참조 구조}: 비용 $c(x,y)$는 두 시각의 응집 장 $u_t, u_s$ \emph{모두}에 의존한다. $\varphi_t(x)$는 $u_t$에, $\varphi_s(y)$는 $u_s$에 의존하므로, 전달 계획 $M_{t \to s}$가 $u_s$에 의존하고, 동시에 $u_s$는 (시간적 에너지를 통해) $M_{t \to s}$에 의존한다. 이 순환적 의존이 자기-참조적 고정점 문제를 야기한다.

\begin{center}
\fbox{\parbox{0.85\textwidth}{\centering
\textbf{자기-참조 순환}\\[4pt]
$u_s$ $\longrightarrow$ $\varphi_s$ $\longrightarrow$ $c(x,y)$ $\longrightarrow$ $M_{t \to s}$ $\longrightarrow$ $\mathcal{E}_{\mathrm{tr}}$ $\longrightarrow$ $u_s$
}}
\end{center}

이 자기-참조 순환을 해결하는 것이 §\ref{sec:fixed-point}의 주제이다.


\section{엔트로피-정칙화 부분 최적 전달}
\label{sec:ot-sinkhorn}

\subsection{Kantorovich 문제}

시각 $t$의 응집 장 $u_t$를 시각 $s$의 $u_s$로 전달하는 문제를 최적 전달(optimal transport)로 정식화한다:
\begin{equation}
    \mathbf{M}^*_{t \to s} = \argmin_{\mathbf{M} \geq 0} \sum_{x,y} c(x,y)\, \mathbf{M}(x,y) + \varepsilon_{\mathrm{OT}} \sum_{x,y} \mathbf{M}(x,y) \log \mathbf{M}(x,y)
    \label{eq:kantorovich}
\end{equation}
subject to:
\begin{align}
    \sum_y \mathbf{M}(x,y) &\leq \mu_t(x) \quad \forall x \quad \text{(원천 제약: 부-확률적)} \label{eq:source-constraint} \\
    \sum_x \mathbf{M}(x,y) &\leq \nu_s(y) \quad \forall y \quad \text{(목표 제약)} \label{eq:target-constraint}
\end{align}

여기서:
\begin{itemize}
    \item $\mu_t(x) = u_t(x)$: 원천 주변 분포(source marginal)
    \item $\nu_s(y) = u_s(y)$: 목표 주변 분포(target marginal)
    \item $\varepsilon_{\mathrm{OT}} > 0$: 엔트로피 정칙화 매개변수
    \item 부등식 제약: \textbf{부분(partial)} 전달---모든 질량이 전달될 필요 없음 (일부는 ``소멸'')
\end{itemize}

\subsection{왜 부분 전달인가}

완전한 전달($\sum_y M(x,y) = \mu_t(x)$)이 아닌 부분 전달을 사용하는 이유:

\begin{enumerate}
    \item \textbf{형성의 탄생과 소멸}: $u_t$에는 있으나 $u_s$에는 없는 형성이 있을 수 있다 (소멸). 역으로 새 형성이 탄생할 수 있다.
    \item \textbf{공리 E1 (부-확률성)}: 시간적 전달 핵은 행 합 $\leq 1$을 만족해야 한다.
    \item \textbf{온톨로지적 일관성}: 응집의 ``보존''을 강제하지 않는다---응집은 약해질 수 있다.
\end{enumerate}

\subsection{로그-영역 Sinkhorn 알고리즘}

엔트로피-정칙화된 최적 전달의 표준적 해법은 Sinkhorn 반복이다. 수치적 안정성을 위해 로그-영역(log-domain)에서 수행한다:

\begin{algorithm}[H]
\caption{로그-영역 Sinkhorn 반복}
\label{alg:sinkhorn}
\begin{enumerate}
    \item \textbf{초기화}: $\mathbf{f} \leftarrow \mathbf{0}$, $\mathbf{g} \leftarrow \mathbf{0}$ (이중 변수)
    \item \textbf{비용 핵}: $\log K(x,y) = -c(x,y) / \varepsilon_{\mathrm{OT}}$
    \item \textbf{반복} ($\ell = 1, 2, \ldots$):
    \begin{enumerate}
        \item $f_i \leftarrow \log \mu_t(i) - \mathrm{LSE}_j(\log K(i,j) + g_j)$ \quad (행 갱신)
        \item $g_j \leftarrow \log \nu_s(j) - \mathrm{LSE}_i(\log K(i,j) + f_i)$ \quad (열 갱신)
        \item 수렴 판정: $\|\mathbf{r} - \mu_t\|_1 < \mathrm{tol}$
    \end{enumerate}
    \item \textbf{출력}: $M(x,y) = \exp(f_x + \log K(x,y) + g_y)$
\end{enumerate}
여기서 $\mathrm{LSE}(\cdot) = \log \sum \exp(\cdot)$는 로그-합-지수 함수이다.
\end{algorithm}

\subsection{수렴 보장}

\begin{theorem}[Sinkhorn 수렴]
\label{thm:sinkhorn-convergence}
$\varepsilon_{\mathrm{OT}} > 0$이고 $\mu_t, \nu_s > 0$이면, 로그-영역 Sinkhorn 반복은 주변 위반(marginal violation) $\ell^1$ 오차가 $O(\log(1/\varepsilon))$ 비율로 0에 수렴하는 유일한 해 $\mathbf{M}^*$에 수렴한다.
\end{theorem}

\textbf{실질적 의미}: $\varepsilon_{\mathrm{OT}}$가 클수록 빠르게 수렴하나 해가 흐려지고(diffuse), 작을수록 정밀하나 수렴이 느려진다. SCC 구현에서는 $\varepsilon_{\mathrm{OT}} = 0.05$--$0.1$이 적절한 균형을 제공한다.

\begin{example}[소규모 Sinkhorn 계산]
\label{ex:small-sinkhorn}
$4 \times 4$ 격자에서의 전달 문제를 수동으로 추적한다. 시각 $t$의 형성이 좌상단에 집중되어 있고 ($\mu = (0.6, 0.3, 0.1, 0)$), 시각 $s$에 약간 이동하여 ($\nu = (0.1, 0.5, 0.3, 0.1)$)라 하자.

비용 행렬 (응집 지문 거리):
\begin{equation}
    c = \begin{pmatrix}
        0.01 & 0.15 & 0.82 & 1.95 \\
        0.18 & 0.02 & 0.21 & 0.98 \\
        0.85 & 0.19 & 0.03 & 0.24 \\
        1.98 & 0.95 & 0.22 & 0.04
    \end{pmatrix} \nonumber
\end{equation}

$\varepsilon_{\mathrm{OT}} = 0.1$로 Sinkhorn 반복을 수행하면:
\begin{itemize}
    \item 반복 1: 주변 오차 $\|r - \mu\|_1 = 0.42$
    \item 반복 5: 주변 오차 $= 0.031$
    \item 반복 10: 주변 오차 $= 0.0018$
    \item 반복 15: 수렴 ($< 10^{-4}$)
\end{itemize}

최종 전달 계획의 구조: $M(1,1) = 0.08$ (지점 1 → 지점 1), $M(1,2) = 0.48$ (지점 1 → 지점 2)가 지배적이며, $M(1,3), M(1,4) \approx 0$. 이는 형성의 핵심(지점 1-2)이 다음 시각의 핵심(지점 2-3)으로 집중적으로 전달됨을 보여준다.
\end{example}

\subsection{$\varepsilon_{\mathrm{OT}}$의 선택과 trade-off}

엔트로피 정칙화 매개변수 $\varepsilon_{\mathrm{OT}}$의 선택은 세 가지 트레이드오프를 반영한다:

\begin{center}
\begin{tabular}{lcc}
\hline
\textbf{기준} & $\varepsilon_{\mathrm{OT}}$ \textbf{큰 경우} ($\geq 0.5$) & $\varepsilon_{\mathrm{OT}}$ \textbf{작은 경우} ($\leq 0.01$) \\
\hline
전달 계획 & 확산적 (균등 배분) & 집중적 (거의 결정론적) \\
수렴 속도 & 빠름 (5--10 반복) & 느림 (50--100+ 반복) \\
수치 안정성 & 높음 & 낮음 (언더플로 위험) \\
자기-참조 고정점 & 수축 보장 쉬움 & 수축 조건 엄격 \\
\hline
\end{tabular}
\end{center}

SCC에서 $\varepsilon_{\mathrm{OT}} = 0.05$--$0.1$이 선호되는 이유: 이 범위에서 전달 계획이 핵심-핵심 전달에 충분히 집중되면서도(2층 집중 유지), 자기-참조 고정점의 Banach 수축이 보장된다.

% [Figure placeholder]
\begin{figure}[t]
\centering
\fbox{\parbox{0.85\textwidth}{\centering\vspace{5cm}
\textbf{[그림 10.1]} 전달 핵 $M_{t \to s}$의 시각화.\\
좌: 비용 행렬 $c(x,y)$ (밝은 색 = 높은 비용). 중: Sinkhorn 수렴 과정 (주변 오차 감소).\\
우: 최종 전달 계획 (밝은 색 = 높은 전달량). 핵심-핵심 연결이 지배적.
\vspace{0.5cm}}}
\caption{Sinkhorn 알고리즘에 의한 전달 핵 계산. 자기-참조적 비용에 의해 응집적 구조가 유사한 지점 간 전달이 집중된다.}
\label{fig:sinkhorn-convergence}
\end{figure}


\section{자기-참조적 고정점}
\label{sec:fixed-point}

\subsection{고정점 방정식}

§\ref{sec:cohesion-fingerprint}에서 지적한 자기-참조 순환을 해결해야 한다. 전달 핵은 다음 고정점 방정식의 해이다:
\begin{equation}
    \mathbf{M}^* = \Phi(\mathbf{M}^*), \quad \text{여기서 } \Phi(\mathbf{M}) = \mathrm{Sinkhorn}\big(c[\varphi_t, \varphi_s(\mathbf{M})],\, \mu_t,\, \nu_s(\mathbf{M})\big)
    \label{eq:fixed-point}
\end{equation}
즉, 전달 핵 $\mathbf{M}$이 결정하는 $u_s$로부터 계산되는 비용으로 다시 전달 핵을 구했을 때, 같은 $\mathbf{M}$이 나와야 한다.

\subsection{Schauder 고정점 정리에 의한 존재}

\begin{theorem}[전달 핵 존재]
\label{thm:transport-existence}
$\Phi$는 컴팩트 볼록 집합 $\{\mathbf{M} \geq 0 : \sum_y M(x,y) \leq \mu_t(x)\}$ 위의 연속 사상이다. Schauder 고정점 정리에 의해 고정점 $\mathbf{M}^*$이 존재한다.
\end{theorem}

존재는 보장되나, 유일성은 추가 조건이 필요하다.

\subsection{수축에 의한 유일성}

약한 체제(weak regime)에서, $\Phi$는 수축 사상이 된다:

\begin{theorem}[Banach 수축 유일성]
\label{thm:contraction-uniqueness}
전달 속박(transport confinement) 조건 하에서 $\Phi$의 Lipschitz 상수가 1 미만이며, Banach 수축 사상 정리에 의해 고정점은 유일하다. 반복
\begin{equation}
    \mathbf{M}^{(\ell+1)} = \Phi(\mathbf{M}^{(\ell)})
    \label{eq:banach-iteration}
\end{equation}
은 기하급수적으로 수렴한다.
\end{theorem}

\subsection{전달 속박}

\begin{definition}[전달 속박]
\label{def:transport-confinement}
전달 계획 $\mathbf{M}$이 \emph{전달 속박}을 만족한다 함은, 핵심-핵심(core-to-core) 전달량이 전체 전달의 대부분을 차지하는 것이다:
\begin{equation}
    \frac{\sum_{x \in \mathrm{Core}_t,\, y \in \mathrm{Core}_s} \mathbf{M}(x,y)}{\sum_{x,y} \mathbf{M}(x,y)} \geq 1 - \delta_{\mathrm{conf}}
    \label{eq:confinement}
\end{equation}
여기서 $\delta_{\mathrm{conf}}$는 작은 양수이다.
\end{definition}

전달 속박의 물리적 의미: 응집 지문의 차이($\Delta_\varphi^2$)가 핵심과 비핵심 사이에서 크게 벌어지므로, 최적 전달은 자연스럽게 핵심-핵심 전달에 집중한다. 구체적으로:
\begin{itemize}
    \item 깊은 핵심 ($\delta \geq 2$): $\Delta_\varphi^2 \approx 2.38$ (큰 지문 간극 → 핵심-비핵심 전달 억제)
    \item 경계 지점: $\Delta_\varphi^2 \approx 0.05$ (작은 지문 간극 → 전달 가능)
\end{itemize}

% [Figure placeholder]
\begin{figure}[t]
\centering
\fbox{\parbox{0.85\textwidth}{\centering\vspace{5cm}
\textbf{[그림 10.2]} 자기-참조적 고정점 반복의 수렴.\\
가로축: 반복 횟수 $\ell$. 세로축: $\|\mathbf{M}^{(\ell+1)} - \mathbf{M}^{(\ell)}\|$.\\
약한 체제에서 기하급수적 수렴. 고정점은 3--5회 반복으로 충분히 근사.
\vspace{0.5cm}}}
\caption{Banach 반복에 의한 자기-참조적 전달 핵의 수렴. 약한 체제에서 기하급수적 수렴이 보장된다.}
\label{fig:fixed-point-convergence}
\end{figure}


\section{T-Persist-1: 단일 형성 지속성}
\label{sec:t-persist-1}

\subsection{왜 지속성 증명이 필요한가}

형성이 시각 $t$에 존재한다는 것은 T1(존재 정리)에 의해 보장된다. 그러나 시각 $t+\Delta t$에도 ``같은'' 형성이 존재한다는 것은 별도의 증명이 필요하다. T-Persist-1은 단일 형성의 지속성을 보장하는 충분 조건을 제시한다.

\subsection{다섯 가지 조건 (a)--(e)}

\begin{theorem}[T-Persist-1]
\label{thm:t-persist-1}
단일 형성 $u_t$가 다음 조건을 만족하면, 시각 $t + \Delta t$에 $u_t$의 구조적 후계자 $u_s$가 존재한다:

\begin{enumerate}
    \item[\textbf{(a)}] \textbf{지지}: $u_t$가 연결 그래프 $X$ 위에 지지됨 (임의의 연결 그래프 허용)

    \item[\textbf{(b)}] \textbf{분지 반경}: 에너지 최소점의 분지 반경 $r \geq 0.210$ (분기로부터 충분히 떨어짐)

    \item[\textbf{(c)}] \textbf{깊은 핵심}: $|\mathrm{Core}| \geq 25$ (등둘레 부등식을 통해 H2'이 보장됨)

    \item[\textbf{(d)}] \textbf{약한 체제}: 전달 고정점이 존재하고 Banach 수축이 적용됨

    \item[\textbf{(e)}] \textbf{전달 속박}: $\mu_0 \geq 3.4$ (깊은 핵심의 지문 간극이 충분히 큼)
\end{enumerate}
\end{theorem}

각 조건의 역할을 살펴보자.

\subsection{조건 (b): 분지 반경과 분기로부터의 거리}

분지 반경 $r$은 에너지 곡면의 극소점 주변에서 그래디언트 흐름이 해당 극소점으로 수렴하는 영역의 ``크기''를 측정한다. $r \geq 0.210$은 시간적 섭동 하에서도 형성이 같은 분지 내에 머무를 만큼 극소점이 ``충분히 깊다''는 것을 보장한다.

이 하한은 방향별 분석(directional analysis)에 의해 도출되었다: 등방적(isotropic) 한계의 2.5--4.3배에 해당하며, 타원체형 분지 구조를 반영한다.

\subsection{조건 (c): 깊은 핵심의 크기와 등둘레 부등식}

형성의 핵심이 충분히 커야 구조적 계승이 의미를 가진다. $|\mathrm{Core}| \geq 25$는 그래프 등둘레 부등식(isoperimetric inequality)을 통해 핵심의 ``내부 깊이''가 보장되는 문턱이다. 이보다 작은 핵심은 본질적으로 ``전이 대역만으로 이루어진'' 형성이며, 구조적 계승의 대상이 되기에 너무 얕다.

\subsection{2층 집중 (Two-Tier Concentration)}

T-Persist-1의 핵심 결론은 전달의 \textbf{2층 집중}이다:

\begin{center}
\fbox{\parbox{0.85\textwidth}{
\textbf{2층 집중 (Two-Tier Concentration)}\\[6pt]
\begin{itemize}
    \item \textbf{깊은 핵심} ($\delta \geq 2$): 전달 질량의 $> 99.99\%$가 핵심-핵심으로 집중
    \item \textbf{얕은 내부} ($\delta = 1$): 전달 질량의 $> 99.3\%$가 내부-내부로 집중
\end{itemize}
여기서 $\delta$는 핵심 경계로부터의 그래프 거리이다.
}}
\end{center}

직관적으로, 깊은 핵심의 지문 간극($\Delta_\varphi^2 \approx 2.38$)이 경계의 지문 간극($\Delta_\varphi^2 \approx 0.05$)보다 압도적으로 크므로, 최적 전달은 핵심의 질량을 핵심으로, 내부의 질량을 내부로 보내는 것을 강하게 선호한다.

해석: 지속성은 ``\emph{질량이 어디로 가는가}''에 관한 것이지, ``\emph{질량이 보존되는가}''에 관한 것이 아니다. 핵심의 질량이 핵심으로 계승되면, 형성은 지속한다---설사 총 질량이 약간 감소하더라도.

\begin{example}[2층 집중의 수치적 검증]
\label{ex:two-tier-numerical}
$15 \times 15$ 격자에서의 단일 형성 지속성 실험. 형성의 핵심은 $(7,7)$ 근방에 위치하며, $|\mathrm{Core}| = 36$ (조건 (c) 만족).

전달 핵 $M_{t \to s}$의 질량 분포:
\begin{center}
\begin{tabular}{lcc}
\hline
\textbf{전달 경로} & \textbf{질량 비율} & \textbf{이론적 한계} \\
\hline
깊은 핵심 → 깊은 핵심 ($\delta \geq 2$) & $99.993\%$ & $> 99.99\%$ \\
얕은 핵심 → 얕은 핵심 ($\delta = 1$) & $99.41\%$ & $> 99.3\%$ \\
경계 → 경계 & $87.2\%$ & --- (이론적 한계 없음) \\
핵심 → 경계 (``누출'') & $0.006\%$ & $< 0.01\%$ \\
핵심 → 외부 (``소멸'') & $0.001\%$ & $< 0.01\%$ \\
\hline
\end{tabular}
\end{center}

주목할 점: 깊은 핵심의 질량은 99.993\%가 다음 시각의 깊은 핵심으로 전달된다. 이는 지문 간극 $\Delta_\varphi^2 \approx 2.38$에 의한 것으로, 핵심에서 경계로의 전달 비용이 핵심-핵심 비용의 약 1700배이기 때문이다.
\end{example}

\subsection{조건 (d)와 (e)의 상호 관계}

조건 (d)(약한 체제)와 (e)(전달 속박)는 독립적이지만 상호 보강한다:

\begin{itemize}
    \item 조건 (d)는 자기-참조 고정점의 \emph{존재와 유일성}을 보장한다. 이것이 없으면 전달 핵 자체가 잘 정의되지 않을 수 있다.
    \item 조건 (e)는 그 고정점의 \emph{구조적 품질}을 보장한다. 전달 핵이 존재하더라도, 핵심-핵심 집중이 없으면 구조적 계승이 의미를 잃는다.
    \item 실질적으로, 조건 (e)의 $\mu_0 \geq 3.4$는 조건 (d)의 수축 조건보다 약간 더 강하다. 따라서 (e)가 만족되면 (d)도 자동으로 만족되는 경우가 대부분이다.
\end{itemize}

\subsection{T-Persist-1의 한계와 열린 문제}

T-Persist-1은 다음과 같은 한계를 가진다:

\begin{enumerate}
    \item \textbf{조건 (b)의 분지 반경 $r \geq 0.210$}: 이 하한은 방향별 분석에서 도출되었으나, 최적(tight)인지는 알려져 있지 않다. 더 작은 $r$에서도 지속성이 가능할 수 있으며, 이는 열린 문제이다.
    \item \textbf{조건 (c)의 $|\mathrm{Core}| \geq 25$}: 작은 그래프에서는 이 조건이 제한적이다. $10 \times 10$ 격자에서 $m = 0.15$인 형성은 핵심이 15 지점 정도이므로 조건을 위반할 수 있다.
    \item \textbf{조건 (d)와 (e)는 카테고리 C}: 이 조건들의 증명은 조건부(conditional)이며, 완전한 증명은 이론의 미래 과제이다.
\end{enumerate}

% [Figure placeholder]
\begin{figure}[t]
\centering
\fbox{\parbox{0.85\textwidth}{\centering\vspace{5cm}
\textbf{[그림 10.3]} 핵심-핵심 계승 (core-to-core inheritance).\\
좌: 시각 $t$의 형성 (핵심=짙은 색, 경계=연한 색). 우: 시각 $s$의 형성.\\
화살표: 전달 핵 $M_{t \to s}$에 의한 질량 흐름. 핵심-핵심 전달이 $> 99.99\%$.
\vspace{0.5cm}}}
\caption{2층 집중에 의한 핵심-핵심 계승. 깊은 핵심의 질량은 거의 전부가 다음 시각의 핵심으로 전달된다.}
\label{fig:core-to-core}
\end{figure}


\section{다중 형성 지속성}
\label{sec:multi-persist}

\subsection{T-Persist-K-Sep: 잘 분리된 체제}

잘 분리된 체제($\Lambda_{\mathrm{coupling}} < 0.01$)에서는 다중 형성 지속성이 단일 형성 지속성으로 환원된다:

\begin{theorem}[T-Persist-K-Sep]
\label{thm:persist-k-sep}
$K$개 형성이 잘 분리되어 있을 때 ($\Lambda < 0.01$), 각 형성 $u^k$가 T-Persist-1의 조건을 독립적으로 만족하면, 모든 형성이 시간적으로 지속한다. 구체적으로:
\begin{enumerate}
    \item[(a)] 형성별 최소점 안정성: $\|u^k_s - u^k_t\| \leq 2\varepsilon_1/\mu_k$
    \item[(b)] 분리 보존: $d_{\min}(j,k)(s) \geq d_{\min}(j,k)(t) - 4\varepsilon_1/\min(\mu_j, \mu_k)$
    \item[(c)] 핵심 포함: $\mathrm{Core}_t(u^k) \subseteq \mathrm{Core}_s(u^k, \theta_{\mathrm{core}} - 2\varepsilon_1/\mu_k)$
    \item[(d)] 전달 집중: 각 $M^k$가 독립적으로 2층 집중을 만족
    \item[(e)] 단체 보존: $\sum_k u^k_s \leq 1 + O(K \exp(-c_0 D_{\mathrm{sep}}))$
\end{enumerate}
\end{theorem}

핵심 아이디어: 잘 분리된 형성들은 서로를 ``보지 못하므로'', 각각의 지속성을 독립적으로 증명할 수 있다.

\subsection{T-Persist-K-Weak: 약하게 상호작용하는 체제}

약하게 상호작용하는 체제($0.01 \leq \Lambda < 1/(K-1)$)에서는 형성 간 간섭이 존재하므로, 추가 조건이 필요하다:

\begin{theorem}[T-Persist-K-Weak]
\label{thm:persist-k-weak}
약하게 상호작용하는 $K$개 형성이 T-Persist-1 조건을 만족하고, 추가적으로 \textbf{격리 조건(isolation condition)}을 만족하면, 모든 형성이 시간적으로 지속한다. 격리 조건: 각 형성의 전달 계획이 다른 형성의 핵심으로의 ``누출''이 제한됨.
\end{theorem}

\subsection{T-Persist-K-Unified: 통합 정리}

제9장의 T-Persist-K-Unified(정리~\ref{thm:persist-k-unified})는 위 두 정리를 $\Lambda_{\mathrm{coupling}}$ 기반의 하나의 매개변수적 정리로 통합한다. 체제가 $\Lambda$에 의해 자동으로 분류되므로, 사용자는 $\Lambda$만 계산하면 적절한 지속성 보장을 얻는다.

\subsection{시간 영역의 단체 제약}

다중 형성의 시간적 전달에서 추가적으로 고려해야 할 제약이 있다:
\begin{equation}
    \sum_{k=1}^{K} M^k_{t \to s}(x, y) \leq M_{\mathrm{total}}(x, y) \quad \forall x, y
    \label{eq:temporal-simplex}
\end{equation}
즉, 하나의 지점에서 출발하는 전달 질량이 여러 형성에 분배되더라도, 총합은 원래의 전달 용량을 초과할 수 없다. 이는 공간 영역의 단체 제약(식~\eqref{eq:simplex-constraint})의 시간적 대응물이다.

% [Figure placeholder]
\begin{figure}[t]
\centering
\fbox{\parbox{0.85\textwidth}{\centering\vspace{5cm}
\textbf{[그림 10.4]} 다중 형성의 시간적 전달.\\
좌: 시각 $t$의 두 형성 (적, 청). 우: 시각 $s$의 두 형성.\\
각 형성의 전달 계획은 독립적이되, 시간적 단체 제약을 만족.\\
잘 분리된 체제에서 교차 전달은 지수적으로 억제됨.
\vspace{0.5cm}}}
\caption{다중 형성의 시간적 전달. 잘 분리된 체제에서 각 형성의 전달 계획은 사실상 독립적이다.}
\label{fig:multi-temporal-transport}
\end{figure}


\section{요약}
\label{sec:transport-summary}

이 장에서 확립된 핵심 사항:

\begin{itemize}
    \item \textbf{추적 ID 없는 정체성}: 시간적 정체성은 추적 라벨이 아닌 응집적 구조의 계승에 의해 성립한다 (발생적 정체성).

    \item \textbf{응집 지문}: 3성분 벡터 $\varphi = (u, \mathrm{Cl}(u), D(\cdot; 1-u))$가 각 지점의 응집적 상태를 요약하며, 자기-참조적 비용을 정의한다.

    \item \textbf{최적 전달}: 엔트로피-정칙화된 부분 최적 전달(로그-영역 Sinkhorn)이 전달 핵을 계산한다.

    \item \textbf{자기-참조 고정점}: Schauder 정리(존재) + Banach 수축(유일성, 약한 체제)이 자기-참조 순환을 해결한다.

    \item \textbf{T-Persist-1}: 5가지 조건 하에서 단일 형성의 지속성이 보장되며, 핵심-핵심 전달은 2층 집중을 보인다 ($>99.99\%$).

    \item \textbf{다중 형성 지속성}: T-Persist-K-Sep(잘 분리된), T-Persist-K-Weak(약하게 상호작용), T-Persist-K-Unified(통합)가 체제별 지속성을 보장한다.
\end{itemize}

다음 장에서는 지금까지 개별적으로 등장한 네 가지 구조적 측면---결합(Bind), 분리(Sep), 내부성(Inside), 지속(Persist)---을 하나의 진단 벡터로 통합하고, 이 벡터가 어떤 이론적 예측을 생성하며, 그 예측들이 어떻게 검증되었는지를 다룬다.
